\documentclass[11pt,a4paper]{amsart}
\usepackage[utf8]{inputenc}
\usepackage[T1]{fontenc}
\usepackage[ngerman]{babel}
\usepackage{amsmath,amssymb,amsthm,mathrsfs}
\usepackage{array,booktabs,longtable}
\usepackage{hyperref}

\newtheorem{theorem}{Satz}[section]
\newtheorem{proposition}[theorem]{Proposition}
\newtheorem{corollary}[theorem]{Korollar}
\newtheorem{definition}[theorem]{Definition}
\newtheorem{remark}[theorem]{Bemerkung}
\newtheorem{openproblem}[theorem]{Offenes Problem}

\newcommand{\stM}{\ensuremath{\checkmark[M]}}
\newcommand{\stKM}{\ensuremath{\checkmark[K/M]}}
\newcommand{\stMneg}{\ensuremath{\checkmark[M]_{\rm neg}}}
\newcommand{\stMpart}{\ensuremath{\checkmark[M]_{\rm part}}}
\newcommand{\Open}{\texttt{?[O]}}
\newcommand{\Conditional}{\textsc{Conditional}}

\title[P05 --- Relative Prime Channels]{P05 --- Relative Prime Channels and Arithmetic Edge Geometry}
\author{Objekt-X-Programm}
\date{2026-08-09}

\begin{document}

\begin{abstract}
Wir konsolidieren die lokale Primkanal- und Primkantengeometrie des
Objekt-X-Programms in typisierter Form. Strikt getrennt werden die Rohkopplung
$T_p$, der von einer gewählten Primhebung induzierte Kanaloperator
$C_p^{[\widehat\varepsilon_p]}$ und seine relative Rang-eins-Realisierung
$C_p^{\rm rel}[\widehat\varepsilon_p]$. Die exakte Liftzulässigkeit ist
nichtlinear, ein allgemeiner exakt zulässiger Nichtnullzeuge ist nicht
konstruiert, und die Herkunft eines geladenen Repräsentanten $L_3^\circ$
bleibt offen. Gesichert sind die kontrollierte modale Rohkopplung, die
feste-$p$-Kollisions- und Restklassenstruktur, die Transportnatur von
$D_{\rm rel}$ in den auditierten Primsektoren, die Spektralmaßdarstellung,
mögliche nichtorthogonale Primkanalbilder und die Trägertrennung direkter
Kreuzprimkollisionen vom von-Mangoldt-Träger. Die vollständige operatorische
Primzahlpotenzrealisierung bleibt konditional.
\end{abstract}

\maketitle

\begin{remark}[Provenienz und Status]\label{rem:provenance}
Kanonische Markdown-Quelle ist
\texttt{papers/P05\_Relative\_Prime\_Channels\_and\_Arithmetic\_Edge\_Geometry.md},
Status \texttt{SYN FINAL AUDITED}, Commit \texttt{bc49413f}. Die Pass-A-Basis
ist Gruppe F, \texttt{PASS A COMPLETE}, Commit \texttt{9c23fc49}. Der
P05-SYN-Primärcheck liegt in
\texttt{audits/AUDIT-2026-08-09\_P05\_SYN\_Primaercheck.md}
(Commit \texttt{54374bdf}), der unabhängige pfadgebundene Zweitcheck in
\texttt{audits/AUDIT-2026-08-09\_P05\_SYN\_Zweitcheck\_Pfadgebunden.md}
(Commit \texttt{f50f3502}) mit dem Urteil
\emph{P05-SYN-GEGENCHECK OHNE KONKRETEN GEGENBEFUND}. Die Synthese hebt keinen
offenen mathematischen Punkt hoch.
\end{remark}

\section{Typisierte Primkanalarchitektur}

\begin{definition}[Drei verschiedene Operatorrollen]\label{def:three-roles}
Für jede Primzahl $p$ werden die drei Rollen strikt getrennt:
\[
\boxed{T_p\neq C_p^{[\widehat\varepsilon_p]}\neq
C_p^{\rm rel}[\widehat\varepsilon_p].}
\]
$T_p$ ist die Rohkopplung auf dem kontrollierten Fourier-/BC-Quellbereich,
$C_p^{[\widehat\varepsilon_p]}$ der von einer gewählten Primhebung induzierte
Kanaloperator und $C_p^{\rm rel}[\widehat\varepsilon_p]$ die relative,
modellrelativ eindimensional induzierte Realisierung. Rang- oder Normaussagen
über $C_p^{\rm rel}$ dürfen nicht auf $T_p$ übertragen werden.
\end{definition}

\begin{remark}[Status]
Die Typtrennung ist die verbindliche F1/F2-Firewall, Status \stKM.
\end{remark}

\begin{theorem}[Nullmodusobstruktion]\label{thm:null-mode}
Auf der kontrollierten modalen Rohkopplungsformel gilt
\[
T_p(e_0V_p)=0.
\]
Der ungechargte Fouriermodus erzeugt daher in diesem lokalen Modell keine
Rohkopplung.
\end{theorem}

\begin{remark}[Status]
Theorem~\ref{thm:null-mode} hat Status \stM\ auf dem kontrollierten
Definitionsbereich; Provenienz NEU-041/155.
\end{remark}

\begin{remark}[SYN-Notation: Rang-eins-Projektion]
Historisch wurde für die orthogonale Rang-eins-Projektion teilweise $\pi_p$
verwendet. Da NEU-158/160 $\pi_p$ zugleich für eine Symmetriedarstellung
benutzen, schreibt P05
\[
\boxed{\Pi_p^{(1)}:=|e_1^{(p)}\rangle\langle e_1^{(p)}|.}
\]
\end{remark}

\begin{theorem}[Rang-eins-Aussage nur für die induzierte Ebene]
\label{thm:rank-one}
Im eindimensional induzierten relativen Modell gilt
\[
\operatorname{rank}C_p^{\rm rel}[\widehat\varepsilon_p]\le1.
\]
Für
\[
P_p=|c_p|^2\Pi_p^{(1)},\qquad
(\Pi_p^{(1)})^2=\Pi_p^{(1)}
\]
folgt
\[
P_p^2=|c_p|^2P_p.
\]
Insbesondere ist $P_p$ im Allgemeinen kein orthogonaler Projektor.
\end{theorem}

\begin{remark}[Status]
Theorem~\ref{thm:rank-one} ist modellrelativ \stM.
\end{remark}

\begin{openproblem}[Intrinsisches Primgewicht]\label{op:intrinsic-cp}
Aus der Rang-eins-Struktur folgt weder unbedingtes $c_p\neq0$ noch
Hebungsunabhängigkeit von $|c_p|^2$ oder die termweise Asymptotik
\[
|c_p|^2\asymp\frac{(\log p)^2}{p}.
\]
Nichtentartung und Hebungsunabhängigkeit bleiben \Open.
\end{openproblem}

\section{Liftgeometrie und exakte Zulässigkeit}

\begin{remark}[SYN-Notation: Liftkern und Liftform]
Die historischen F2-Blätter verwenden $K_p$ für den Kern der primitiven
Projektion. Da F3 $K_p,K_{pq}$ für Prim-/Feshbachblöcke verwendet, setzt P05
\[
\boxed{\mathscr K_p^{\rm lift}:=\ker\pi_{\rm prim}.}
\]
Die verbundene Liftform wird zur Abgrenzung vom späteren BC-Halbgewicht als
\[
\boxed{h_p^{\rm conn}}
\]
geschrieben. Beide Änderungen sind reine SYN-Disambiguierungen.
\end{remark}

\begin{definition}[Primitive Liftfaser]
Die Fourier-geladenen Richtungen
\[
E_p^{\rm ch}:=\operatorname{span}\{e_uV_p:u\neq0\}
\]
liegen algebraisch in $\mathscr K_p^{\rm lift}$.
\end{definition}

\begin{theorem}[Quadratische Normierungsbedingung]\label{thm:quadratic-lift}
Für eine Liftänderung $k$ relativ zu einem Basispunkt
$\widehat\varepsilon_p^{\,0}$ besitzt die exakte Normierungsbedingung die Form
\[
2\operatorname{Re}h_p^{\rm conn}(\widehat\varepsilon_p^{\,0},k)
+h_p^{\rm conn}(k,k)=0.
\]
Sie ist daher keine homogene lineare Kerngleichung.
\end{theorem}

\begin{remark}[Status]
Theorem~\ref{thm:quadratic-lift} hat Status \stM\ (NEU-157 rev.3 / 165b).
\end{remark}

\begin{theorem}[Keine zusätzliche lineare Kernfamilie im auditierten Quellenkegel]
\label{thm:no-extra-linear-kernel}
Im explizit auditierten Quellenkegel liefern NEU-157 und NEU-44 keine
zusätzlichen nichttrivialen Operatoren
\[
L_{p,a}:\mathscr K_p^{\rm lift}\longrightarrow Y_{p,a}
\]
mit exakter Zulässigkeitsbedingung $L_{p,a}(k)=0$. In diesem Quellenkegel gilt
\[
\mathscr K_p^{\rm hom}=\mathscr K_p^{\rm lift}.
\]
\end{theorem}

\begin{remark}[Scope]
Theorem~\ref{thm:no-extra-linear-kernel} ist \stMneg\ als Quellenbefund,
nicht als globaler mathematischer Unmöglichkeitssatz.
\end{remark}

\begin{openproblem}[Exakt zulässiger Nichtnullzeuge]
Ein vollständiger Beweis der Existenz eines exakt zulässigen $k$ mit
\[
T_p(k)\neq0
\]
liegt im auditierten Quellenstand nicht vor. Präprojektive und infinitesimale
Satzschemata sind gültig; die Integration in die exakte nichtlineare
Zulässigkeitsmenge bleibt \Open.
\end{openproblem}

\section{Rohkopplung und feste-$p$-Kollisionsgeometrie}

\begin{definition}[Modale Rohkopplung]\label{def:modal-raw}
Für $u\neq0$ lautet die kontrollierte präprojektive Formel
\[
T_p^{\rm pre}(e_uV_p)
=-\sum_{s,m}\ell_{s,m}\,u\,s\log p\;e_{u+ps}V_{pm}.
\]
Sie ist ausschließlich auf dem kontrollierten modalen Quellbereich zu
verwenden und liefert keine automatische globale Operatorverlängerung.
\end{definition}

\begin{theorem}[Injektivität bei festem Eingangsmodus]
Für festes $p$ und festes $u$ ist
\[
(s,m)\longmapsto(u+ps,pm)
\]
injektiv. Verschiedene Summanden eines einzelnen Eingangsmodus kollidieren
nicht in derselben algebraischen Zielkoordinate. Status \stM.
\end{theorem}

\begin{corollary}[Bedingte Einzelmoden-Nichtverschwindung]
Definiere
\[
\operatorname{supp}^{\times}(L_3^\circ)
:=\{(s,m):s\ell_{s,m}\neq0\}.
\]
Dann gilt auf dem kontrollierten modalen Bereich
\[
\operatorname{supp}^{\times}(L_3^\circ)\neq\varnothing
\Longrightarrow T_p^{\rm pre}(e_uV_p)\neq0\qquad(u\neq0).
\]
Die Nichtleerheitsvoraussetzung ist im alten Quellenpfad nicht hergeleitet;
die Herkunft bleibt \Open.
\end{corollary}

\begin{theorem}[Restklassen-/Faltungszerlegung]\label{thm:convolution}
Für eine Zielkollision innerhalb eines festen Primkanals gilt
\[
(u,s,m)\sim_p(u',s',m')
\Longleftrightarrow
m=m',\qquad u-u'=p(s'-s).
\]
Nach Zerlegung in Restklassen $c\in\mathbb Z/p\mathbb Z$ wird das
Mehrmoden-Kernproblem zu
\[
x^{(c)}*b^{(m)}=0.
\]
\end{theorem}

\begin{remark}[Firewall]
Theorem~\ref{thm:convolution} betrifft einen festen Primkanal $p$ und ist
von einer Kreuzprimkollision $p\neq q$ der Form $pm_p=qm_q$ zu trennen.
\end{remark}

\section{Herkunftsblockade des geladenen $L_3^\circ$}

\begin{proposition}[Rechenzulässig ist nicht herkunftszulässig]
Die Modellwahl $L_3^\circ=e_1V_1$ ist innerhalb des historischen Testmodells
rechenzulässig und liefert etwa $(p-1)\log p\neq0$. Daraus folgt nicht, dass
$e_1V_1$ aus $[L_3]$ hergeleitet oder kanonisch gewählt ist. Die Herkunft ist
\Conditional/\Open.
\end{proposition}

\begin{theorem}[Typ- und Ladungsblockade sind unabhängig]
NEU-170d trennt
\[
[L_3]\not\longmapsto L_3^\circ=e_1V_1
\]
und
\[
(p-1)\log p\neq0
\not\Longrightarrow C_p(e_{1-p}V_p)\neq0.
\]
Ein positiver Skalarfaktor konstruiert weder einen typisierten Repräsentanten
noch die benötigte nichtverschwindende relative Zielkante.
\end{theorem}

\begin{proposition}[Endstatus des alten Typquellenpfads]
Im auditierten Quellenkegel ist kein vollständiges Tupel
\[
(B_3,M,C^\bullet,b,L_3,\rho_{\rm op})
\]
konstruiert. Dies ist \stMneg\ als Quellenbefund und \Open\ als mathematische
Neukonstruktionsfrage. Außerdem ist $\delta_{\rm BC}:A_Q\to A_Q$ eine
Algebraableitung und nicht das Hochschild-Kodifferential
$b:C^n(B_3,M)\to C^{n+1}(B_3,M)$.
\end{proposition}

\section{Quotienten- und Symmetriearchitektur}

\begin{remark}[SYN-Notation: Quotient und Darstellung]
NEU-159/160 verwenden $Q_p$ für den Rohkopplungsquotienten. Zur Abgrenzung
vom lokalen Weil-Beitrag schreibt P05
\[
\boxed{\mathscr Q_p^{\rm quot}:=Q_p^{(\mathrm{NEU\text{-}159/160)}}}.
\]
Die historische Symmetriedarstellung wird als
\[
\boxed{\pi_p^{\rm sym}:G_p\to\mathcal U(\mathscr Q_p^{\rm quot})}
\]
geschrieben.
\end{remark}

\begin{proposition}[Abstrakte Quotientenstruktur]
Gesichert bleiben die positive Quotientenform nach Nullraumfaktorisierung,
die skaliert-isometrische Identifikation mit dem abgeschlossenen Bildraum,
das allgemeine Nullraumabstiegslemma und das Intertwining-zu-Unitärität-Lemma.
Status \stM.
\end{proposition}

\begin{theorem}[Abstrakter Kommutantensatz]
Für eine konkrete unitäre Darstellung $\pi_p^{\rm sym}$ auf
$\mathscr Q_p^{\rm quot}$ sind beschränkte positive semidefinite
$G_p$-invariante Formen genau dann skalare Vielfache einer Referenzform, wenn
\[
\pi_p^{\rm sym}(G_p)'=\mathbb C I.
\]
Status \stM.
\end{theorem}

\begin{openproblem}[Konkrete Quotientenrealisierung]
Nicht automatisch mitbewiesen sind
$\mathscr Q_p^{\rm quot}\neq\{0\}$, eine konkrete unitäre Wirkung
$\pi_p^{\rm sym}$, deren Irreduzibilität und die konkrete Eindeutigkeit der
verbundenen Form. Status \Open/\Conditional.
\end{openproblem}

\begin{remark}[Firewall]
NEU-166a/166b liefern keine unbedingte globale Erweiterung von $T_p^{\rm pre}$,
keinen vollständigen Definitionsbereich, keinen Quotientendeszent und keinen
kanonischen transversalen Detektor. Fall 3a ist nur lokal/modenweise
formelmäßig bestätigt; die globale Fallentscheidung bleibt \Open.
\end{remark}

\section{Transportgeometrie der auditierten Primsektoren}

\begin{remark}[SYN-Notation: Transportkoeffizient]
NEU-225 verwendet historisch ebenfalls $c_p$ für den Transportkoeffizienten.
P05 disambiguiert rein redaktionell durch
\[
\boxed{\kappa_p^{\rm tr}:=\frac12\gamma_Np\log p.}
\]
\end{remark}

\begin{theorem}[Transportnormalform]\label{thm:transport}
Auf einer auditierten Primfaser $\mathcal H_{p,a}$ gilt
\[
D_{\rm rel}\big|_{\mathcal H_{p,a}}
\cong2i\kappa_p^{\rm tr}\frac{d}{dt}
\quad\text{auf}\quad
L^2(\mathbb R)\oplus L^2(\mathbb R).
\]
Status \stM.
\end{theorem}

\begin{corollary}[Spektraltyp im auditierten Primsektor]
In den auditierten Primsektoren besitzt $D_{\rm rel}$ rein absolutstetiges
Spektrum und keinen Kern; ein kompakter Resolvent ist ausgeschlossen. Damit
ist $D_{\rm rel}$ dort ein Transportgenerator und kein bereits gefundener
Hilbert--Pólya-Operator.
\end{corollary}

\begin{openproblem}[Zusammengesetzte Sektoren]
Für nichtprimitive beziehungsweise zusammengesetzte $m$-Sektoren können
Mehrfachsprünge Restklassen mischen. Eine globale Spektralaussage über alle
Mischsektoren ist nicht bewiesen. \texttt{[O-225-3]} bleibt \Open.
\end{openproblem}

\begin{theorem}[Spektralmaß statt diskreter Eigenbasis]\label{thm:spectral-measure}
Die historische diskrete Eigenbasisdarstellung wird durch die
projektionswertige Spektralmaßform ersetzt. Für $V_p,V_q$ definiere
\[
\mu_{pq}^{a,b}(B):=\langle V_pa,E_D(B)V_qb\rangle.
\]
Dann lautet die verbindliche Resolventenschreibweise
\[
\langle a,K_{pq}(z)b\rangle
=\int_{\mathbb R}\frac{d\mu_{pq}^{a,b}(\lambda)}{\lambda-z}.
\]
Status \stKM.
\end{theorem}

\section{Nichtorthogonale Primkanalgeometrie}

\begin{definition}[Koordinatenwörterbuch]
\[
\eta_{p;m;s,u}\longleftrightarrow e_RV_M,
\qquad M=pm,\quad R=u+ps.
\]
Diese Koordinatenidentifikation ist insbesondere im Primsektor vollständig
kontrolliert. Status \stM.
\end{definition}

\begin{theorem}[Primkanalbilder müssen nicht orthogonal sein]
\label{thm:nonorthogonal}
Verschiedene Primkanalbilder können nichttrivial überlappen. Daher ist
\[
\mathcal K_N\stackrel{?}{=}\bigoplus_pK_p
\]
nicht strukturell erzwungen. Kreuzblöcke $K_{pq}$ können generisch
nichtverschwinden; dies behauptet nicht $K_{pq}\neq0$ für jedes $p\neq q$.
Status \stM.
\end{theorem}

\begin{remark}[Mechanismus]
Die Off-Diagonalität kann aus
$\operatorname{Ran}V_p\not\perp\operatorname{Ran}V_q$ entstehen, ohne dass
$D_{\rm rel}$ selbst Primlabels mischen muss.
\end{remark}

\section{Arithmetische Primzahlpotenzgewichte}

\begin{proposition}[Primitiver algebraischer Halbgewichtsfaktor]
Im primitiven $p$-Kanal wird algebraisch
\[
h_p^{\rm bal}=p^{-1/2}I
\]
und damit der lokale Faktor $\log p/\sqrt p$ erhalten. Status \stMpart;
Hilbert-Selbstadjungiertheit, Abschließbarkeit, Domäne und globaler
Funktionalkalkül sind damit nicht bewiesen.
\end{proposition}

\begin{theorem}[Testfunktionswert als Matrixkoeffizient]
Für die Autokorrelation gilt
\[
g_a(x)=\operatorname{Re}\langle a,U_xa\rangle_{L^2(\mathbb R)},
\]
insbesondere
\[
\boxed{g_a(\log p)=\operatorname{Re}\langle a,U_{\log p}a\rangle.}
\]
Der Ausdruck ist ein unitärer Matrixkoeffizient und kein Normquadrat. Status
\stM.
\end{theorem}

\begin{theorem}[Arithmetische Primzahlpotenzidentität]
\label{thm:mangoldt-primepower}
Für jede Primzahlpotenz $p^m$ gilt
\[
\boxed{\frac{\Lambda(p^m)}{\sqrt{p^m}}
=\frac{\log p}{p^{m/2}}.}
\]
Dies ist eine arithmetische Identität, Status \stM.
\end{theorem}

\begin{remark}[Firewall: keine vollständige Operatorrealisierung]
Im Quellenstand ist kein vollständiger Beweis
\[
h_n^{\rm bal}=n^{-1/2}I\qquad(n\ge1)
\]
gefunden. Die Rückreferenz in NEU-250i ist durch NEU-250g nicht gedeckt.
Daher bleibt
\[
h_{p^m}^{\rm bal}\bigl(H_{\rm pr}^{1/2}E_R,H_{\rm pr}^{1/2}E_{R'}\bigr)
=\frac{\log p}{p^{m/2}}\delta_{RR'}
\]
\Conditional.
\end{remark}

\begin{remark}[Firewall: $H_{\rm pr}$ ist nicht $\Lambda$]
Die formal gradnormierte BC-Energie darf nicht mit $\Lambda$ auf allen
natürlichen Zahlen identifiziert werden. Beispielsweise gilt $\Lambda(6)=0$,
während die gradnormierte Energie nicht automatisch verschwindet.
\end{remark}

\section{Trägertrennung: direkte Kreuzprimkollision versus Mangoldt-Träger}

\begin{theorem}[Kreuzprimkollision erzeugt keinen Mangoldt-Trägerpunkt]
\label{thm:support-separation}
Seien $p\neq q$ Primzahlen und $pm_p=qm_q=M$. Dann besitzt $M$ mindestens
zwei verschiedene Primteiler und daher $\Lambda(M)=0$. Folglich
\[
\boxed{\operatorname{supp}\Lambda\cap
\operatorname{supp}(\text{direkte Kreuzprimkollision})=\varnothing.}
\]
Status \stM.
\end{theorem}

\begin{remark}[Firewall]
Theorem~\ref{thm:support-separation} widerspricht
Theorem~\ref{thm:nonorthogonal} nicht. Die Trägertrennung beweist weder
globale Primorthogonalität noch globales Verschwinden der Kreuzblöcke.
\end{remark}

\section{Schnittstellen zu den Folgesynthesen}

\begin{proposition}[Routing nach P06]
Nach P06 gehen Schur-/Feshbach-Komplementfragen, globale Quotienten- und
Deszentfragen, Spektralmaß-/Schattenklassenkriterien, $u$-Regulator und
Quellhilbertisierung sowie die offenen zusammengesetzten Sektoren
\texttt{[O-225-3]}.
\end{proposition}

\begin{proposition}[Routing nach P09]
Nach P09 gehen das Typfundament
\[
(B_3,M,C^\bullet,b,L_3,\rho_{\rm op}),
\]
die Konstruktion eines typkorrekten Repräsentanten von $[L_3]$ und die
strikte Trennung von $\delta_{\rm BC}$ und Hochschild-$b$.
\end{proposition}

\begin{proposition}[Routing nach P11]
Nach P11 gehen gemeinsame globale Quellenräume, Gramoperatoren der
überlappenden Primkanalbilder, globale Off-Diagonalblöcke, der Mediatorstatus
J-A/J-B und die Verbindung zwischen arithmetischer Trägerstruktur und globaler
Weil-Geometrie. P05 behauptet keine dieser globalen Konstruktionen bereits
gelöst zu haben.
\end{proposition}

\section{Statusmatrix}

\small
\begin{longtable}{>{\raggedright\arraybackslash}p{0.51\textwidth}
>{\raggedright\arraybackslash}p{0.23\textwidth}
>{\raggedright\arraybackslash}p{0.16\textwidth}}
\toprule
Aussage & Status & Quelle/Paket\\
\midrule
\endfirsthead
\toprule
Aussage & Status & Quelle/Paket\\
\midrule
\endhead
$T_p\neq C_p^{[\widehat\varepsilon_p]}\neq C_p^{\rm rel}[\widehat\varepsilon_p]$
& \stKM\ Typfirewall & F1/F2\\
$T_p(e_0V_p)=0$ & \stM & F1/F2\\
$\operatorname{rank}C_p^{\rm rel}\le1$ im induzierten Modell
& \stM\ modellrelativ & F1/F2\\
intrinsisches $c_p\neq0$ / Nichtentartung & \Open & F1/F2\\
Hebungsunabhängigkeit & \Open & F2\\
quadratische Normierungsbedingung & \stM & F2\\
zusätzliche lineare Kernfamilie im auditierten Quellenkegel
& \stMneg\ Quellenbefund & F2\\
exakt zulässiger Nichtnullzeuge & \Open & F2\\
feste-$p$-Injektivität und Faltungskollision & \stM & F2\\
geladener $L_3^\circ$ aus $[L_3]$ hergeleitet & \Open & NEU-170d/173\\
abstrakte Quotienten-/Intertwining-Lemmata & \stM & F2\\
konkrete Quotientenrealisierung / Irreduzibilität
& \Open/\Conditional & F2\\
Transportnormalform in Primsektoren & \stM & F3\\
reines a.c.-Spektrum / kein Kern in auditierten Primsektoren & \stM & F3\\
globale Spektralaussage in zusammengesetzten Sektoren
& \Open & \texttt{[O-225-3]}\\
projektionswertige Spektralmaßform & \stKM & F3\\
Primkanalbilder können nichttrivial überlappen & \stM & F3\\
Primblockdiagonalität ist nicht strukturell erzwungen & \stM & F3\\
primitiver Faktor $\log p/\sqrt p$ & \stMpart & F4\\
$g_a(\log p)=\operatorname{Re}\langle a,U_{\log p}a\rangle$ & \stM & F4\\
$\Lambda(p^m)/\sqrt{p^m}=\log p/p^{m/2}$ & \stM & F4\\
$h_n^{\rm bal}=n^{-1/2}I$ für alle $n$
& \Conditional; Beweis nicht gefunden & F4\\
vollständige operatorische Primzahlpotenzrealisierung & \Conditional & F4\\
Kreuzprimkollision $\cap\operatorname{supp}\Lambda=\varnothing$ & \stM & F4\\
Trägertrennung impliziert keine Primorthogonalität & \stM & F3/F4\\
\bottomrule
\end{longtable}
\normalsize

\section{Offene Kernfragen nach P05}

\begin{enumerate}
\item \textbf{Liftintrinsik:} hebungsunabhängige, nichtentartete Primkanalgröße
$|c_p|^2$.
\item \textbf{Exakte Zulässigkeit:} exakt zulässiger geladener Lift mit
nichtverschwindender Rohkopplung.
\item \textbf{$L_3$-Herkunft:} typisierte Repräsentantenbrücke von $[L_3]$.
\item \textbf{Globale Deszendenz:} vollständiger Definitionsbereich und
Quotientendeszent der Rohkopplung.
\item \textbf{Zusammengesetzte Sektoren:} Spektral- und Restklassengeometrie
jenseits der auditierten Primfasern.
\item \textbf{Primzahlpotenzoperator:} allgemeine BC-Halbgewichtung und
Hilbert-Fundierung des gradnormierten Operators.
\item \textbf{Globale Gramgeometrie:} nichtorthogonale globale Kopplung der
lokalen Primkanalbilder mit dem archimedischen Kanal.
\end{enumerate}

\section{SYN-Endurteil}

Der konsolidierte lokale Bestand
\[
\boxed{\begin{gathered}
\text{lokale Rohkopplung}+\text{nichtlineare Liftgeometrie}\\
+\text{Transportprimfasern}+\text{nichtorthogonale Kanalbilder}\\
+\text{arithmetische Trägertrennung}
\end{gathered}}
\]
ist miteinander kompatibel, bildet aber noch keine globale
Objekt-X-Realisierung. Der Schutzsatz der weiteren Migration lautet
\[
\boxed{\begin{gathered}
\text{lokale Primkanalstruktur}\neq\text{globale positive Gramkopplung}\\
\neq\text{Hilbert--Pólya-Operator}
\end{gathered}}.
\]

\begin{remark}[SYN-Endstatus]
Die Markdown-Inhaltsstufe ist \texttt{SYN FINAL AUDITED}. Diese LaTeX-Fassung
enthält keine Statushochstellung und keine neue mathematische Behauptung. Der
LaTeX-SYN-Transferaudit ergab keinen Status-, Typ-, Formel- oder
Routingkonflikt. Endstatus: \texttt{SYN FROZEN} \stKM.
\end{remark}

\end{document}
